\documentclass[11pt]{article}
\usepackage[a4paper,margin=27mm]{geometry}
\usepackage[T1]{fontenc}
\usepackage{lmodern,amsmath,amssymb,amsthm,booktabs,microtype}
\usepackage[colorlinks=true,allcolors=blue]{hyperref}
\newtheorem{theorem}{Theorem}
\newtheorem{lemma}[theorem]{Lemma}
\newtheorem{proposition}[theorem]{Proposition}
\newtheorem{corollary}[theorem]{Corollary}
\newcommand{\Tr}{\operatorname{Tr}}
\newcommand{\id}{\operatorname{id}}
\newcommand{\supp}{\operatorname{supp}}
\newcommand{\CEA}{C_{\mathrm{EA}}}
\newcommand{\Rcal}{\mathcal R}
\newcommand{\Tcal}{\mathcal T}
\newcommand{\comp}{\mathrm{cmp}}
\title{Sharp entropy budgets for communication\\with unknown resets and optical loss}
\author{Lluis Eriksson}
\date{26 October 2026}
\begin{document}
\maketitle
\begin{abstract}
A memoryless reset channel transmits the input with fixed probability
\(0<\lambda<1\) and otherwise prepares a state \(\tau\).
We investigate recovery of its flagged-erasure capacity benchmarks when
only an entropy budget on the reset is available, with no bound on its
energy or orientation. For oscillator input energy \(E\), and each fixed
finite R\'enyi order \(\alpha>0\), the sharp conditions for uniform absolute
entanglement-assisted recovery are \(2^{B(E)}/E\to0\) for \(\alpha<1\),
\(B(E)/\log_2(E+1)\to0\) for \(\alpha=1\), and \(B(E)\to0\) for
\(\alpha>1\). The sufficient conditions recover four benchmarks even with
one encoder and decoder for the entire compound channel class.
For Shannon and fixed super-Shannon order we further determine the full
bounded entanglement-assisted minimax loss, uniformly over every budget:
\(b(\lambda,q\min\{1,B/\log_2(E+1)\})\) and
\(b(\lambda,q[1-2^{-((\alpha-1)/\alpha)B}])\), respectively, up to
absolute \(o(1)\), where \(q=1-\lambda\) and
\(b(u,v)=(u+v)h(u/(u+v))\).
Known-noise and unknown-fixed-noise minimax losses coincide asymptotically.
The same loss curves hold with a quantum-limited attenuator in the
transmitted branch, uniformly over transmissivities bounded away from
zero, relative to its attenuated flagged benchmark. One common assisted
code handles an unknown fixed reset and an unknown fixed transmissivity.
A fixed finite output compression, spectral branch-information bounds and
exact type-subspace blocking preserve the physical input constraint.
A vacuum--thermal mixture supplies the converse through a finite-flag
continuity argument. Finite-codimension energy avoidance also gives uniform
pointwise recovery of two-way benchmarks. Spectral approximation estimates
and compound coding theorems are established ingredients; the claimed
addition concerns their operational entropy-budget quantifiers and the
matching uniform loss curves.
The noise is fixed throughout a code block. No error-exponent, arbitrarily
varying-channel or universal two-way coding result is asserted.
\end{abstract}

\section{Model and the distinction between two kinds of uniformity}
Let \(\mathcal H=\ell^2(\mathbb N_0)\), let
\(H=\sum_{j\ge0}j|j\rangle\langle j|\), and put \(q=1-\lambda\), with
\(0<\lambda<1\) fixed. The channel is
\[
 \Rcal_\tau(\rho)=\lambda\rho+q\tau\Tr\rho .
\]
The reset is a density operator on \(\mathcal H\). Its eigenvectors need
not have finite energy. Write
\[
 g(E)=(E+1)\log_2(E+1)-E\log_2E,\qquad
 h(x)=-x\log_2x-(1-x)\log_2(1-x).
\]
All information quantities are in bits. The six capacity conventions and
their erasure coefficients are
\[
\begin{array}{c|cccccc}
 X&C&\CEA&Q&P&Q_2&K_2\\ \hline
 r_X&\lambda&2\lambda&(2\lambda-1)_+&(2\lambda-1)_+&\lambda&\lambda .
\end{array}                                                    \tag{1}
\]
The input constraint is expected block-average physical energy at most
\(E\). Entanglement assistance is unlimited in \(\CEA\).
For \(Q_2,K_2\), local operations are countably decomposable LOCC with
free authenticated public discussion and no initial entanglement or key;
the energy accounting includes source copies consumed to create a seed.
Security includes a full channel environment and the public transcript.
These are the conventions of \cite{prior,dsw}. Local output preparation
is not charged as channel input energy.

For a class \(\Tcal_E\) define the pointwise uniform deficit
\[
 \Delta_X(E,\Tcal_E)=
 \sup_{\tau\in\Tcal_E}\{r_Xg(E)-X(\Rcal_\tau,H,E)\}.              \tag{2}
\]
Here the code may depend on \(\tau\).
In contrast, \(X_\comp(E,\Tcal_E)\), for \(X=C,\CEA,Q,P\), denotes the
compound capacity: both users know the class, neither knows \(\tau\),
and one coding sequence must have error tending to zero uniformly for
all its members. The same \(\tau\) is used on every channel use.
Classical error uses the average-message convention; private communication
requires trace secrecy for a uniform message against the full environment,
uniformly in the reset. Trace secrecy means one half of the trace norm
between the actual message--environment state and the product of its
uniform-message marginal with its actual average environment marginal.
The construction below also permits vanishing
unnormalized Holevo leakage. \(Q\) uses entanglement transmission.
Compound capacities are bounded
above by every individual capacity. Reset is receiver postprocessing of
flagged erasure, so all quantities are at most \(r_Xg(E)\)
\cite{prior,dsw}. In particular \(Q=P=0\) if \(\lambda\le1/2\).
We always take the coding limit at fixed \(E\) first.

For fixed finite \(\alpha>0\), let \(S_\alpha\) be the R\'enyi entropy,
with \(S_1=S\), and set
\[
 \Tcal_{\alpha,B}(E)=\{\tau:S_\alpha(\tau)\le B(E)\},\qquad
 S_\alpha(\tau)=\frac{\log_2\Tr\tau^\alpha}{1-\alpha}
 \quad(\alpha\ne1).                                           \tag{3}
\]
The budget is any finite nonnegative function. The class is nonempty
and includes arbitrary pure states.

\begin{theorem}[Entropy-budget classification]\label{thm:main}
For fixed \(0<\lambda<1\) and fixed finite \(\alpha>0\), the following
are equivalent as \(E\to\infty\):
\begin{enumerate}
\item The corresponding condition in this table holds:
\[
\begin{array}{c|c}
 \text{order}&\text{condition}\\ \hline
 0<\alpha<1&2^{B(E)}/E\longrightarrow0\\
 \alpha=1&B(E)/\log_2(E+1)\longrightarrow0\\
 1<\alpha<\infty&B(E)\longrightarrow0 .
\end{array}                                                    \tag{4}
\]
\item \(\Delta_{\CEA}(E,\Tcal_{\alpha,B}(E))\to0\).
\item \(\CEA^{\comp}(E,\Tcal_{\alpha,B}(E))=2\lambda g(E)+o(1)\).
\item All four compound capacities \(C,\CEA,Q,P\) equal their respective
benchmarks \(r_Xg(E)+o(1)\).
\end{enumerate}
These conditions also imply \(\Delta_X(E,\Tcal_{\alpha,B}(E))\to0\)
for all six entries of (1).
\end{theorem}

Necessity is asserted through \(\CEA\), and thus for simultaneous
recovery; it is not asserted separately for every unassisted capacity.
No uniformity near \(\alpha=1\), near an endpoint of \(\lambda\), or in
the blocklength required at changing \(E\) is claimed.

The finite-dimensional identity/replacement family appears explicitly in
\cite[Eq.~(34)]{glr}; its output-norm multiplicativity question differs
from the constrained capacity problem here. The erasure comparison is also
a prior ingredient. Our earlier work \cite{prior} supplied pointwise spectral
recovery and a finite-state distillation estimate.
The energy-budget minimax problem in \cite{minimax} constrains a different
reset class. Its phase-twirling reduction cannot simply be used here:
dephasing an equal-amplitude pure superposition on a finite support
raises every finite-order R\'enyi entropy
from zero to the logarithm of its support size.
R\'enyi spectral-tail control is an established approximation tool
\cite{vc}, and qualitative differences among entropy orders already
occur in matrix-product approximation \cite{schuch}.
The compound coding theorems used below are likewise prior results
\cite{cq,quantumcompound,eacompound}.
The vacuum-replacement special case is an established binary fading
model, with a finite-energy input optimization in
\cite[Section IV and Appendix C]{catalano}. Thermal inputs need not be
optimal there. Our capacity converses do not assume that they are.
The present statements do not claim those ingredients or a generic
entropy-order distinction as discoveries.
Energy-constrained compound cq theory also exists \cite[Theorem 5]{cnr}.
Its uniform output-energy and effective-projection hypotheses do not hold
for our full reset classes: for any finite-rank output projection a pure
reset can lie orthogonally to it, leaving reset mass \(q\) outside.
Our fixed output processing is not claimed to approximate every physical
output in trace norm; it directly preserves the information bounds needed
for coding.

\section{A fixed input and output compression}
Write the decreasing eigenvalues of \(\tau\) as \(p_j(\tau)\), indexed
from one, and define
\[
 \eta_m(\Tcal)=\sup_{\tau\in\Tcal}
       \left(1-\sum_{j=1}^m p_j(\tau)\right).
\]
A spectral-tail condition does not imply trace-norm compactness when
all eigenvector orientations are allowed.

Fix \(E>2\), put \(E'=E-1\), \(r=E'/(E'+1)\), and choose
\[
 L=\lceil4E\ln(E+1)\rceil,\quad \epsilon=r^L,\quad
 P_L=\sum_{j=0}^{L-1}|j\rangle\langle j|.
\]
Let \(\sigma\) be the thermal state of mean \(E'\), truncated and
normalized on \(P_L\). It is independent of the reset. Direct summation gives
\begin{align}
 \Tr H\sigma&=E'-\frac{L\epsilon}{1-\epsilon}<E,\notag\\
 \|\sigma\|_\infty&=\frac{1}{E(1-\epsilon)},\notag\\
 s_E:=S(\sigma)&=g(E')+\log_2(1-\epsilon)
       -\frac{L\epsilon}{1-\epsilon}\log_2(1/r).                 \tag{5}
\end{align}
Since \(\epsilon\le(E+1)^{-4}\), \(g(E)-s_E\to0\).

The fixed receiver channel
\[
 D_L(X)=P_LXP_L\ \oplus\ \Tr[(I-P_L)X]|f\rangle\langle f|
\]
has output dimension \(L+1\).
Identify the logical input with \(P_L\mathcal H\); then the decoded reset
channel is
\[
 N_\tau(\rho)=\lambda\rho+q\nu_\tau\Tr\rho,\qquad
 \nu_\tau=D_L(\tau).                                          \tag{6}
\]
The physical and logical transmit state in (6) agree on the first block.

\begin{lemma}[Compression and branch information]\label{lem:branch}
For each integer \(m\ge1\), the spectral tail of \(\nu_\tau\) after rank
\(m+1\) is at most the tail of \(\tau\) after rank \(m\).
Set
\[
 z_E=\frac{\lambda(m+1)}{E(1-\epsilon)}
                    +q\eta_m(\Tcal_E).                        \tag{7}
\]
If \(z_E\le1/2\), uniformly for \(\tau\in\Tcal_E\),
\begin{align}
 I(\sigma,N_\tau)&\ge2\lambda s_E-h(z_E),\notag\\
 I_c(\sigma,N_\tau)&\ge(2\lambda-1)s_E-h(z_E),\notag\\
 \chi(\{p_j,N_\tau(|j\rangle\langle j|)\})&
                 \ge\lambda s_E-h(z_E),                       \tag{8}
\end{align}
where \(p_j\) are the diagonal probabilities of \(\sigma\).
\end{lemma}
\begin{proof}
Let \(\tau_m\) be a top-\(m\) spectral truncation. Then
\(0\le D_L(\tau_m)\le\nu_\tau\), its trace is at least
\(1-\eta_m\), and its rank is at most \(m+1\). Its support therefore
witnesses the claimed tail bound by the variational characterization
of the largest eigenvalue sum.

Introduce a classical transmit/reset register \(F\), of probabilities
\(\lambda,q\). A hypothetical measurement onto a subspace of rank at most \(m+1\)
carrying its largest spectral mass guesses reset on success. Its error is at most (7):
the transmit contribution is bounded by
\(\lambda(m+1)\|\sigma\|_\infty\), and the reset contribution by \(q\eta_m\).
Binary Fano and measurement data processing give \(H(F|B)\le h(z_E)\).

For a purification of \(\sigma\), \(I(R:BF)=2\lambda s_E\);
the reference is independent of \(F\).
The loss on forgetting the flag is
\(I(R:F|B)\le H(F|B)\), since \(F\) is classical and
\(H(F|RB)\ge0\). Subtracting \(S(R)=s_E\) gives the coherent-information
bound. For the classical Fock ensemble the same chain rule starts from
\(I(X:BF)=\lambda s_E\). All systems here are finite dimensional.
\end{proof}
The hypothetical test depends on \(\tau\) only to bound information.
It is not the implemented decoder. The actual common decoder comes
from compound coding after the fixed \(D_L\).

\section{Compound coding with an exact physical energy constraint}
An unconstrained finite-dimensional capacity formula by itself does not
prove a physical input-cost statement. The following block construction
supplies that step.

\begin{lemma}[Exact type-subspace blocking]\label{lem:type}
Let \(p\) be a rational distribution on finitely many energy eigenvectors,
with mean energy strictly below \(E\), and let \(\sigma_p\) be its diagonal
state. At blocklength \(k\) with integral counts \(kp_j\), let \(T\)
be its exact type subspace and \(\pi_T\) its uniform state. Put
\[
 d_k=kH(p)-\log_2\dim T=O(\log k)
\]
with the input alphabet fixed. Every vector in \(T\) has total energy
\(k\sum_jjp_j<kE\). For every finite-dimensional channel \(N\),
\begin{align}
 I(\pi_T,N^{\otimes k})&\ge kI(\sigma_p,N)-2d_k,\notag\\
 I_c(\pi_T,N^{\otimes k})&\ge kI_c(\sigma_p,N)-d_k.              \tag{9}
\end{align}
The uniform ensemble of Fock strings in \(T\) has Holevo information
at least \(k\chi(p,N)-d_k\).
\end{lemma}
\begin{proof}
Each one-site marginal of \(\pi_T\) is \(\sigma_p\), and
\(D(\pi_T\Vert\sigma_p^{\otimes k})=d_k\).
Relative-entropy data processing bounds output total correlation by
\(d_k\); hence the output entropy is at least
\(kS(N(\sigma_p))-d_k\).
The complementary-output entropy is at most
\(kS(N^c(\sigma_p))\) by subadditivity.
Together with \(S(\pi_T)=kH(p)-d_k\), these prove (9).
A finite Stinespring environment suffices.
For the classical assertion, the average output is \(N^{\otimes k}(\pi_T)\),
while the averaged conditional output entropy is exactly
\(k\sum_jp_jS(N(|j\rangle\langle j|))\).
The energy assertion holds throughout the type subspace, including
superpositions.
\end{proof}

\begin{proposition}[Uniform compound sufficient criterion]\label{prop:compound}
If integers \(m(E)\ge1\) satisfy
\[
 m(E)/E\longrightarrow0,\qquad
 \eta_{m(E)}(\Tcal_E)\longrightarrow0,                          \tag{10}
\]
then all four compound benchmarks in Theorem~\ref{thm:main} are recovered
with absolute error tending to zero.
For \(z_E\le1/2\), a sufficient deficit bound is
\[
 r_X[g(E)-s_E]+h(z_E),\qquad X=C,\CEA,Q,P,                      \tag{11}
\]
with the zero \(Q,P\) cases treated separately.
\end{proposition}
\begin{proof}
At fixed \(E,L\), approximate the distribution in (5) by rational
distributions with strict mean-cost slack.
All information quantities converge uniformly over the decoded class:
input and output dimensions are fixed, and finite-dimensional entropy
continuity has constants independent of the channel.
For each rational distribution apply Lemma~\ref{lem:type}.

For \(C\), use the compound cq coding theorem
\cite[Theorems III.1 and V.10]{cq} on the alphabet of exact-type Fock
strings, with its uniform distribution.
For \(Q\), the finite-dimensional compound result
\cite[Section 6.1, Lemma 16]{quantumcompound} applies directly to the
maximally mixed subspace input \(\pi_T\), giving its infimum coherent
information as an achievable rate with one encoder and recovery.
For \(\CEA\), use
\cite[Theorem 1 and Lemma 9]{eacompound}, together with its equivalence
between entanglement-assisted classical and quantum transmission.
These theorems allow arbitrary index sets and uninformed users.

Every encoder of the restricted block channel takes values in tensor
powers of \(T\). The energy constraint is therefore automatic for every
physical input, not merely for its average over random code constructions.
Divide rates by \(k\), let the type blocklength grow to remove \(d_k/k\),
and then take the rational approximation limit.
Equations (8) give (11).
For privacy, let \(\varepsilon_n\) be the worst entanglement infidelity
of the common quantum code. Uhlmann's theorem and contraction of trace
distance, followed by measurement of the reference in the message basis,
give uniform-message trace secrecy at most \(2\sqrt{\varepsilon_n}\)
relative to the actual average environment state. The decoder's basis
measurement has average success at least its entanglement fidelity.
Include the physical environment and all systems discarded by \(D_L\)
and the decoder in that environment. No extra assistance is used.
For unnormalized Holevo secrecy one can choose net radius \(n^{-6}\)
in the proof of \cite[Lemma 16 and Eq.~(34)]{quantumcompound}, at fixed
\(E,p,k\) and fixed rate backoff. The net size is polynomial; its coding
term decays exponentially, and tensor-product transfer costs at most
\(n\,n^{-6}\). Thus \(\varepsilon_n=O(n^{-5})\).
Since the message alphabet size \(M_n\) has \(\log M_n=O(n)\),
finite-message conditional-entropy continuity \cite[Lemma 2]{winter}
also makes the Holevo leakage vanish. This does not assert a uniform
error exponent as \(E\) grows. No shared randomness or channel knowledge
is imported into \(C,Q,P\).
Finally (5), (7), and (10) make the right side vanish.
\end{proof}
For any fixed rate slack, choose finite type blocklength and rational
distribution first, and then take the outer coding blocklength to infinity.
Taking suprema over those finite choices gives the displayed rate bound.
The finite-block coding theorems are used as existing results, not
reproved here. The type-space estimates address the extra energy issue.
They do not control the resulting blocklength or reliability exponent.
The distinction matters: universal error exponents need not equal
channel-aware exponents \cite{reliability}.

\section{Entropy budgets give the sufficient spectral condition}
For \(0<\alpha<1\), ordering the eigenvalues and putting
\(R=\sum_jp_j^\alpha\) gives \(p_j\le(R/j)^{1/\alpha}\).
Summation bounded by an integral yields
\[
 \eta_m\le
 \min\left\{1,\frac{\alpha}{1-\alpha}
       \left(\frac{2^B}{m}\right)^{(1-\alpha)/\alpha}\right\}.    \tag{12}
\]
R\'enyi control of spectral approximation tails is prior work
\cite[Lemma 2]{vc}; (12) is a convenient nonsharp elementary form.
If \(2^{B(E)}/E\to0\), choose
\(m=\lfloor\sqrt{E\,2^{B(E)}}\rfloor\). Then both parts of (10) hold.

For Shannon entropy, \(p_j\le1/j\) gives
\[
 \eta_m\le\min\{1,B/\log_2(m+1)\}.                              \tag{13}
\]
Choose \(m=\lfloor\sqrt E\rfloor\) under the middle condition of (4).
For \(\alpha>1\),
\[
 \sum_jp_j^\alpha\le p_1^{\alpha-1},\qquad
 \eta_1\le1-2^{-B}.                                           \tag{14}
\]
Thus \(m=1\) suffices when \(B\to0\).
These proofs are for each fixed order, with constants allowed to depend
on it. They establish the compound sufficiency in Theorem~\ref{thm:main}.

\section{A converse that excludes improved input states}
It is not enough to show a positive loss for one input.
We show that a vacuum--thermal reset also prevents all near-entropy-optimal
inputs from removing that loss.

For any \(\rho,\tau\), choose a purification \(\psi_\rho\) and define
\[
 \Omega_{JRB}=
 \lambda|0\rangle\langle0|_J\otimes\psi_\rho+
 q|1\rangle\langle1|_J\otimes\rho_R\otimes\tau_B,\qquad
 e_\tau(\rho)=I(J:R|B)_\Omega .
\]
Relative-entropy definitions make this well defined and
\(0\le e_\tau(\rho)\le h(\lambda)\).
For finite input entropy,
\[
 I(\rho,\Rcal_\tau)=2\lambda S(\rho)-e_\tau(\rho).              \tag{15}
\]
This is the established finite-branch identity of \cite{prior}.
For finite-energy input the entropy is finite.

Fix \(0<a<1\), let \(\gamma_E\) be the thermal state of mean \(E\), set
\(p_0=1/(E+1)\), and choose
\[
 \tau_{E,a}=(1-a)|0\rangle\langle0|+a\gamma_E,\qquad
 b_0=(\lambda+qa)h\left(\frac{\lambda}{\lambda+qa}\right)>0.     \tag{16}
\]
At input \(\gamma_E\), diagonal observation outside the vacuum has fixed
branch ratio \(\lambda:qa\). Hence \(H(J|B)\ge(1-p_0)b_0\).
The test onto its purification has false-transmit probability
\(q\Tr(\gamma_E^2\tau_{E,a})\le qp_0^2\), and never misses transmission.
For large \(E\), binary Fano therefore gives
\[
 e_{\tau_{E,a}}(\gamma_E)\ge(1-p_0)b_0-h(qp_0^2)
                          \longrightarrow b_0.               \tag{17}
\]

\begin{lemma}[Uniform finite-branch continuity]\label{lem:continuity}
For input trace distance \(t=\|\rho-\gamma\|_1/2\), uniformly in \(\tau\),
\[
 |e_\tau(\rho)-e_\tau(\gamma)|\le
 \omega(\min\{1,\sqrt{2t}\}),\quad
 \omega(s)=4s+2(1+s)h\left(\frac{s}{1+s}\right).                \tag{18}
\]
No entropy or energy constraint on the conditioning systems is required.
\end{lemma}
\begin{proof}
Choose Uhlmann purifications in a common reference space.
Their trace distance is at most \(\sqrt{2t}\); the reference marginals
are no further apart. The two branch states with the same \(\tau\)
therefore have that distance bound.
Use \(e=H(J|B)-H(J|RB)\) and apply the finite-system
conditional-entropy continuity inequality twice, with \(\dim J=2\)
\cite[Lemma 2]{winter}. This inequality permits arbitrary conditioning
dimension. Conditional entropy here is \(H(J)-I(J:C)\), not a difference
of two infinite entropies.
Explicitly, compress the conditioning system onto increasing finite
subspaces plus a failure flag. Data processing and relative-entropy
lower semicontinuity show that its mutual information with \(J\)
converges to the original, bounded by \(H(J)\).
Apply the finite-dimensional bound and pass to the limit.
\end{proof}

For any admissible input set \(d=g(E)-S(\rho)\).
The Gibbs identity and Pinsker in bits give
\[
 D(\rho\Vert\gamma_E)
 =d+(\Tr H\rho-E)\log_2(1+1/E)\le d,\qquad
 \tfrac12\|\rho-\gamma_E\|_1\le\sqrt{(\ln2)d/2}.                \tag{19}
\]
Consequently (18) is bounded by
\(\omega(\min\{1,(2\ln2\,d)^{1/4}\})\).
Choose \(d_*>0\) so this modulus at \(d_*\) is less than \(b_0/4\).
For all sufficiently large \(E\), (17) is at least \(b_0/2\).
Using (15), every input then has benchmark deficit at least
\[
 \min\{2\lambda d_*,b_0/4\}>0:                                \tag{20}
\]
if \(d\ge d_*\) use its entropy deficit; otherwise use (18).
The constrained entanglement-assisted capacity formula
\cite[Theorem 1]{hs} now bounds the capacity itself below the benchmark
by (20). No optimizer is assumed to exist.

It remains to fit this same witness into each failed budget condition.
For Shannon entropy,
\(S(\tau_{E,a})\le h(a)+ag(E)\).
If \(B/\log_2(E+1)\not\to0\), choose a subsequence
\(B\ge c\log_2(E+1)\) and a fixed \(0<a<\min\{c,1\}\).
The witness is eventually admissible there.
For \(\alpha\ne1\), with \(r=E/(E+1)\), the exact power sum is
\[
 \Tr\tau_{E,a}^\alpha
 =(1-a+ap_0)^\alpha+
       \frac{a^\alpha p_0^\alpha r^\alpha}{1-r^\alpha}.         \tag{21}
\]
For \(\alpha<1\) it gives
\[
 S_\alpha(\tau_{E,a})
 =\log_2E+\frac{\alpha\log_2a-\log_2\alpha}{1-\alpha}+o(1).
\]
If \(2^B/E\not\to0\), a subsequence has \(B\ge\log_2E-C\).
Choose fixed \(a\) small enough that the constant offset is less than
\(-C-1\). For \(\alpha>1\), (21) instead gives
\[
 S_\alpha(\tau_{E,a})\longrightarrow
            -\frac{\alpha}{\alpha-1}\log_2(1-a).
\]
If \(B\not\to0\), take a subsequence \(B\ge c>0\) and make this limit
less than \(c/2\).
Thus (20) obstructs uniform individual \(\CEA\) recovery in all failed
regimes. Compound capacity cannot be larger.
Together with Proposition~\ref{prop:compound}, this proves the four
equivalences in Theorem~\ref{thm:main}.

\section{Pointwise two-way recovery without an energy bound on noise}
For completeness we identify the separate mechanism needed for the
last two entries of (1). It does not supply a common two-way code.

\begin{lemma}[Finite-codimension energy avoidance]\label{lem:minmax}
For a discrete Hamiltonian with eigenvalues
\(e_0\le e_1\le\cdots\), let \(M\) be the orthogonal complement of at
most \(m\) vectors, which may themselves have infinite energy.
Any finite distribution \(p_0,\ldots,p_{L-1}\) can be realized by a
state supported in \(M\), with entropy \(H(p)\) and cost at most
\(\sum_{j=0}^{L-1}p_je_{j+m}\).
\end{lemma}
\begin{proof}
Intersect \(M\) with the span of the first \(L+m\) energy vectors.
Its dimension is at least \(L\). The \(j\)th lowest eigenvalue of the
finite compression of \(H\) is at most \(e_{j+m}\): intersecting
instead with the first \(j+m+1\) vectors supplies a subspace of
dimension at least \(j+1\) for the min-max principle.
Put weights \(p_j\) on its first \(L\) eigenvectors.
All chosen vectors lie in a finite energy window; no projection onto
an infinite-energy vector is used to define an input state.
\end{proof}

Avoid the top \(m\) reset eigenvectors and use finite strict-cost
approximants to the oscillator tail entropy \(g(E-m)\), for \(E>m\).
For fixed \(E,m\), first choose \(m<E\prime<E\) with
\(g(E\prime-m)\) as close to \(g(E-m)\) as desired. Upper truncation
of the shifted geometric distribution lowers its mean and converges in
entropy. The resulting finite probability vector has strict cost slack.
Lemma~\ref{lem:minmax} uses that same vector for every forbidden
subspace, so its cost and entropy errors are uniform over orientations. With physical isometry \(V\), logical cost
\(G=V^*HV\) is a bounded finite matrix, diagonal in the constructed basis.
Projection onto \(\supp V\) plus failure flag accepts reset weight
\(\delta\le q\eta_m\).
The dimension-free finite-state recovery estimate in
\cite[Section 2, theorem ``An arbitrary noisy maximally correlated
component'']{prior} gives
\[
 Q_2,K_2\ge
 [\lambda S(\sigma)-b(\lambda,\delta)-\delta\log_2 5]_+,\quad
 b(\lambda,\delta)=(\lambda+\delta)
                 h\left(\frac{\lambda}{\lambda+\delta}\right). \tag{22}
\]
This is an inherited distillation/coding theorem, not established by
the min-max lemma alone. Its seed comes from copies of the same
strictly admissible finite input; their number and energy are counted.
The decoder's larger environment includes the physical environment,
so its security guarantee retains the public transcript.
The prior proof and the erasure bounds use the LOCC/cost conventions
stated in Section 1, not a bound on energy of local operations.
For the inherited upper bounds, see the proposition on erasure upper
bounds in Section 6 of \cite{prior}, and the lemma ``Erasure with a single
finite partition sum'' in its spectral-criterion section. The common-code
four-capacity argument of Sections 2--3 above does not use the two-way
distillation theorem.

Taking capacity suprema removes finite-approximation errors. Monotonicity
of \(b\) in \(\delta\) gives the uniform deficit bound
\[
 \Delta_X(E,\Tcal_E)\le
 \lambda[g(E)-g(E-m)]+b(\lambda,q\eta_m)+q\eta_m\log_2 5,
 \qquad X=Q_2,K_2.                                           \tag{23}
\]
The choices (12)--(14) have \(m/E\to0\) and \(\eta_m\to0\);
\(g(E)-g(E-m)\le m\log_2(1+1/(E-m))\to0\).
This completes the six-benchmark sufficient assertion.

\section{Full entanglement-assisted minimax losses}\label{sec:curves}
The classification has a sharper form for Shannon and super-Shannon
budgets. Put
\[
 b(u,v)=(u+v)h\!\left(\frac{u}{u+v}\right),\qquad
 f(t)=b(\lambda,qt),\qquad 0\le t\le1,
\]
with continuous boundary values. For a scalar budget \(B\ge0\), write
\[
 \begin{aligned}
 D_{\rm known}(E,B)&=\Delta_{\CEA}(E,\{\tau:S_\alpha(\tau)\le B\}),\\
 D_{\comp}(E,B)&=2\lambda g(E)-\CEA^{\comp}
                  (E,\{\tau:S_\alpha(\tau)\le B\}).
 \end{aligned}
                                                               \tag{24}
\]
The order \(\alpha\) is implicit in these two quantities.
Always \(D_{\rm known}\le D_{\comp}\).

\begin{theorem}[Uniform minimax loss curves]\label{thm:curves}
Fix \(0<\lambda<1\). For \(\alpha=1\), or any fixed finite
\(\alpha>1\), set
\[
 t_\alpha(E,B)=
 \begin{cases}
 \min\{1,B/\log_2(E+1)\},&\alpha=1,\\
 1-2^{-((\alpha-1)/\alpha)B},&\alpha>1.
 \end{cases}                                                   \tag{25}
\]
Then
\[
 \sup_{0\le B<\infty}\max_{Y\in\{{\rm known},\comp\}}
 \left|D_Y(E,B)-f(t_\alpha(E,B))\right|\longrightarrow0.
                                                               \tag{26}
\]
In particular the difference between the known-noise and compound
deficits vanishes uniformly in the budget. Vacuum--thermal mixtures
are uniformly asymptotically least favorable. No finite-energy
optimality or uniqueness of the reset is asserted.
\end{theorem}

The result is not obtained by evaluating a single input and calling its
deficit a capacity converse. The next argument controls all inputs. It
uses the branch functional and spectral estimate of \cite{minimax}; the
new matching concerns the entropy-constrained adversary and common code.

\subsection{A branch estimate without the small-error cutoff}
For the fixed decoded channel in (6), take a top-\((m+1)\) reset
projection as in Lemma~\ref{lem:branch}. Write
\(x=\lambda\Tr(P\sigma)\) and
\(\delta=q\Tr((I-P)\nu_\tau)\). Its binary observation \(Y\) has
\[
 H(F|Y)=b(\lambda-x,\delta)+b(x,q-\delta)
 \le b(\lambda,q\eta_m)+b(\bar x,q),\quad
 \bar x=\min\left\{\lambda,
             \frac{\lambda(m+1)}{E(1-\epsilon)}\right\}.        \tag{27}
\]
Here \(\eta_m\) is the supremum tail of the class. The partial
derivatives of \(b(u,v)\) are \(\log_2(1+v/u)\) and
\(\log_2(1+u/v)\), hence both arguments are increasing, with continuous
extensions at zero. Data processing gives \(H(F|B)\le H(F|Y)\).
The same branch chain rule and type-space compound coding as in
Section~3 therefore imply
\[
 D_{\comp}(E,B)\le b(\lambda,q\eta_m)+R_+(E,m),\qquad
 R_+(E,m)=2\lambda[g(E)-s_E]+b(\bar x,q).                       \tag{28}
\]
For \(m/E\to0\), \(R_+\to0\). The hypothetical observation still only
bounds information; it is not the operational common decoder. Capacity
suprema at each fixed energy remove the rational/type approximation
slacks before taking the energy limit.

\subsection{A uniformly sharp witness against every input}
For \(0\le a\le1\), let \(\tau_{E,a}=(1-a)|0\rangle\langle0|+a\gamma_E\),
and put \(p_0=1/(E+1)\). The flagged receiver output at input \(\gamma_E\)
has conditional entropy
\[
 (1-p_0)b(\lambda,qa)+
 b\bigl(\lambda p_0,q(1-a+ap_0)\bigr)
 \le f(a)+b(\lambda p_0,q).
\]
The branch information lost is no larger. This proves the capacity
lower bound needed for the right side of (29) below.

For arbitrary \(\rho\) of mean at most \(E\), define the inherited functional
\[
 F(\rho,\tau)=\lambda S(\rho)+S(\lambda\rho+q\tau)
                         -qS(\tau)-h(\lambda).
\]
The exact identity is \(I(\rho,\Rcal_\tau)=F(\rho,\tau)+H(F|RB)\),
where the last \(F\) is the classical branch register. All entropies
in this witness argument are finite.
Write \(P=\lambda+qa\) and \(\chi=(\lambda\rho+qa\gamma_E)/P\).
Its mean is at most \(E\). Entropy of mixing and maximal thermal entropy give
\(S(\lambda\rho+q\tau)\le h(q(1-a))+Pg(E)\).
A vacuum measurement of the witness's mixture flag has error \(ap_0\), so
\(S(\tau)\ge h(a)+ag(E)-h(ap_0)\).
Using
\(h(\lambda)+qh(a)=h(q(1-a))+P h(\lambda/P)\) yields
\[
 F(\rho,\tau_{E,a})\le2\lambda g(E)-f(a)+q h(ap_0).
\]

Put \(d=g(E)-S(\rho)\), \(r=\|\rho\|_\infty\).
The Gibbs identity (19), followed by measurement in a largest-eigenvalue
vector, gives
\(r\log_2(E+1)-1\le D(\rho\Vert\gamma_E)\le d\).
The purification-projector test has error
\(q\Tr(\rho^2\tau)\le qr^2\).
Always guessing reset has error \(\lambda\). Choosing the better test and
applying binary Fano bounds the conditional branch entropy by
\(h(\min\{\lambda,qr^2\})\); the displayed argument is at most \(1/2\).

If \(d\ge h(\lambda)/(2\lambda)\), the flagged bound
\(I\le2\lambda S(\rho)\) already implies
\(I\le2\lambda g(E)-h(\lambda)\le2\lambda g(E)-f(a)\).
Otherwise let
\[
 \bar r_E=\min\left\{1,
       \frac{1+h(\lambda)/(2\lambda)}{\log_2(E+1)}\right\},\quad
 J_E=h(\min\{\lambda,q\bar r_E^2\}),\quad
 R_-(E)=q h(p_0)+J_E .
\]
For \(E\ge2\), \(h(ap_0)\le h(p_0)\). The global functional bound and
\(r\le\bar r_E\) consequently give, for all \(a\in[0,1]\),
\[
 f(a)-R_-(E)\le
 2\lambda g(E)-\CEA(\Rcal_{\tau_{E,a}},H,E)
 \le f(a)+b(\lambda p_0,q).                                  \tag{29}
\]
Both errors vanish independently of \(a\). The capacity supremum needs
neither an optimizer nor an interchange of near-optimal-input limits.

\subsection{Matching every budget}
The function \(f\) is increasing, concave and zero at zero; its derivative
on \((0,1]\) is \(q\log_2(1+\lambda/(qt))\).
Integrating the decreasing derivative gives
\( |f(x)-f(y)|\le f(|x-y|)\) for \(x,y\in[0,1]\).
This endpoint modulus avoids an invalid bounded-derivative assumption.

For Shannon entropy set \(A=\log_2(E+1)\), \(G=g(E)\),
\(m=\lfloor E/\ln(E+1)\rfloor\), \(L_m=\log_2(m+1)\),
\(d_+=1-L_m/A\) and \(d_-=(G-A+1)/G\), for \(E\ge2\).
The ordered-spectrum bound is \(\eta_m\le\min\{1,B/L_m\}\).
The admissible witness has
\(a=\min\{1,\max\{0,(B-1)/G\}\}\): below budget 1 use the pure vacuum,
in the middle use \(S(\tau)\le1+aG\), and above \(G+1\) use \(\gamma_E\).
With \(t=\min\{1,B/A\}\), the two clipped ratios differ by at most
\(d_+\), while \(0\le t-a\le d_-\), for every \(B\ge0\).
Thus (28)--(29) imply the explicit budget-independent sandwich
\[
 f(t)-f(d_-)-R_-\le D_{\rm known}\le D_{\comp}
       \le f(t)+f(d_+)+R_+(E,m).                              \tag{30}
\]
All remainders tend to zero.

For fixed \(\alpha>1\), put \(R=2^{(1-\alpha)B}\),
\(m=\lfloor\sqrt E\rfloor\), \(t_m=(m+1)^{1-\alpha}\).
If \(w_m\) is the mass of the first \(m\) eigenvalues, then
\(R\le\sum_jp_j^\alpha\le w_m^\alpha+t_m\).
Hence
\[
 \eta_m\le1-(R-t_m)_+^{1/\alpha},\qquad
 0\le R^{1/\alpha}-(R-t_m)_+^{1/\alpha}\le t_m^{1/\alpha}.
\]
The witness with \(a=1-R^{1/\alpha}\) is exactly admissible because
\(\Tr\tau_{E,a}^\alpha\ge(1-a)^\alpha=R\).
Writing \(d_\alpha=t_m^{1/\alpha}\), the matching sandwich is
\[
 f(a)-R_-\le D_{\rm known}\le D_{\comp}
       \le f(a)+f(d_\alpha)+R_+(E,m).                         \tag{31}
\]
This includes \(R\le t_m\), \(B=0\), and unbounded growing budgets.
General admissible resets may have infinite Shannon entropy and energy;
(28) only uses their finite decoded states. The converse only uses the
finite-energy witness. No subtraction of infinite entropies is involved.
Equations (30)--(31) prove Theorem~\ref{thm:curves}.


\section{An unknown optical attenuation in the transmitted branch}
\label{sec:optics}
Let \(\mathcal A_\kappa\) be the quantum-limited attenuator with photon
transmissivity \(0<\kappa\le1\); its environment is in the vacuum.
Consider
\[
 \mathcal N_{\kappa,\tau}(\rho)
      =\lambda\mathcal A_\kappa(\rho)+q\tau\Tr\rho,\qquad
 c_\kappa(E)=g(E)+g(\kappa E)-g((1-\kappa)E).                  \tag{32}
\]
The channel parameters are fixed throughout each block. Its flagged
entanglement-assisted benchmark is \(\lambda c_\kappa(E)\).
The expression \(c_\kappa\) is the established attenuator assisted
capacity \cite{hw}; the entropy bound below also verifies its optimizer.
The vacuum-reset special case is the erasure-lossy fading channel in
\cite[Appendix F, Eq.~(F2)]{catalano}. We do not claim that model as new.

For the reset class \(\Tcal_{\alpha,B}=\{\tau:S_\alpha(\tau)\le B\}\),
define the fixed-\(\kappa\) deficits
\[
 \begin{aligned}
 D_{\kappa,{\rm known}}(E,B)
   &=\sup_{\tau\in\Tcal_{\alpha,B}}
          [\lambda c_\kappa(E)-\CEA(\mathcal N_{\kappa,\tau},H,E)],\\
 D_{\kappa,\comp}(E,B)
   &=\lambda c_\kappa(E)-\CEA^\comp
             (E,\{\mathcal N_{\kappa,\tau}:\tau\in\Tcal_{\alpha,B}\}).
 \end{aligned}                                                 \tag{33}
\]

\begin{theorem}[Uniform optical-loss extension]\label{thm:optics}
Fix \(0<\lambda<1\), \(0<\kappa_0\le1\), and \(\alpha=1\) or a
fixed finite \(\alpha>1\). With \(f,t_\alpha\) from Theorem~\ref{thm:curves},
\[
 \sup_{\substack{0\le B<\infty\\\kappa_0\le\kappa\le1}}
 \max_{Y\in\{{\rm known},\comp\}}
 |D_{\kappa,Y}(E,B)-f(t_\alpha(E,B))|\longrightarrow0.           \tag{34}
\]
If \(\kappa\) is also unknown and belongs to the entire interval
\([\kappa_0,1]\), one common coding sequence for both unknown parameters
has deficit
\[
 \lambda c_{\kappa_0}(E)-\CEA^\comp
 (E,\{\mathcal N_{\kappa,\tau}:\kappa\in[\kappa_0,1],
                             \tau\in\Tcal_{\alpha,B}\})
       = f(t_\alpha(E,B))+o(1),                               \tag{35}
\]
uniformly in finite \(B\ge0\).
\end{theorem}

This theorem concerns assisted communication only. Its benchmark is not
\(2\lambda g(E)\) when \(\kappa<1\). No thermal environment, uniform
limit as \(\kappa_0\downarrow0\), or varying transmissivity within a
codeword is asserted.

\subsection{A finite-energy entropy difference}
For mean \(\mu\le M\), relative-entropy data processing under
\(\mathcal A_t\), which sends \(\gamma_M\) to \(\gamma_{tM}\), gives
\[
 \begin{aligned}
 S(\rho)-S(\mathcal A_t(\rho))
 &\le \log_2\frac{M+1}{tM+1}
  +\mu\left[\log_2(1+1/M)-t\log_2(1+1/(tM))\right]\\
 &\le g(M)-g(tM),\qquad 0<t\le1 .
 \end{aligned}                                                 \tag{36}
\]
The bracket is nonnegative: \(x\log(1+1/(xM))\) increases in \(x>0\),
as its derivative is \(\log(1+y)-y/(1+y)>0\).
For \(t=0\) the same conclusion is \(S(\rho)\le g(M)\).
This is a Gibbs/data-processing estimate, not an assumption of a Gaussian
input. Its unrestricted entropy-gain analogue is prior work \cite{gain}.

The complementary pure-loss output is unitarily equivalent to
\(\mathcal A_{1-\kappa}(\rho)\). Equation (36) and the output entropy
bound \(S(\mathcal A_\kappa(\rho))\le g(\kappa E)\) give
\(I(\rho,\mathcal A_\kappa)\le c_\kappa(E)\), with equality at
\(\gamma_E\). The constrained assisted coding theorem \cite{hs}
identifies this maximum with capacity.
For \(\kappa\ge1/2\), applying (36) to the degraded output with
transmissivity \((1-\kappa)/\kappa\) also gives
\[
 S(\mathcal A_\kappa(\rho))-S(\mathcal A_{1-\kappa}(\rho))
 \le K_\kappa(E):=[g(\kappa E)-g((1-\kappa)E)]_+ .
\]
For \(\kappa\le1/2\) the upper bound is zero by antidegradability and
weak monotonicity, so the displayed inequality holds throughout.

\subsection{A finite-rank test for a mixed transmitted state}
For a witness let \(\tau_{\kappa,E,a}=(1-a)|0\rangle\langle0|
                                       +a\gamma_{\kappa E}\).
The transmitted reference/receiver state
\(\omega_\rho=(\id\otimes\mathcal A_\kappa)(\psi_\rho)\) is mixed;
the purification-projector argument for an identity signal is unavailable.
The exact branch functional now is
\[
 F_\kappa(\rho,\tau)
 =\lambda[S(\rho)-S(\mathcal A_{1-\kappa}(\rho))]
      +S(\lambda\mathcal A_\kappa(\rho)+q\tau)-qS(\tau)-h(\lambda).
\]
It satisfies \(I(\rho,\mathcal N)=F_\kappa+H(F|RB)\).
Using (36) and the same three-weight entropy argument as in Section~7
gives, for every input,
\[
 F_\kappa(\rho,\tau_{\kappa,E,a})
 \le\lambda c_\kappa(E)-f(a)+q h\!\left(\frac{a}{1+\kappa E}\right).
                                                               \tag{37}
\]

Put \(D_0=\log_2(1/\kappa_0)+h(\lambda)/\lambda\).
Flagging gives \(I(\rho,\mathcal N)\le\lambda[S(\rho)+K_\kappa(E)]\).
Uniformly for \(\kappa\ge\kappa_0\),
\(g(E)-c_\kappa(E)+K_\kappa(E)\le\log_2(1/\kappa_0)\):
above \(1/2\) the left side is zero; below it, use
\(g((1-\kappa)E)-g(\kappa E)\le\log_2((1-\kappa)/\kappa)\).
If \(d=g(E)-S(\rho)>D_0\), the input already loses at least
\(h(\lambda)\ge f(a)\) from the flagged benchmark.
For the remaining inputs the earlier spectral argument gives
\(r=\|\rho\|_\infty\le\bar r_E=\min\{1,(D_0+1)/\log_2(E+1)\}\).

Write pure-loss Kraus vectors as
\(|v_\ell\rangle=(I\otimes K_\ell)|\psi_\rho\rangle\), indexed by
the number \(\ell\) of lost photons. Their mean index is
\((1-\kappa)\Tr H\rho\).
Set \(t_E=\sqrt{\ln(E+1)}\), \(K=\lceil t_E(1-\kappa)E\rceil\).
Keeping \(\ell\le K\) loses mass at most \(1/t_E\) by Markov's bound;
for \(\kappa=1\) only \(\ell=0\) occurs.
First project these vectors onto receiver non-vacuum and let \(Q\) be
the support projection of their span. Its rank is at most \(K+1\).
The signal vacuum probability obeys
\[
 \Tr[|0\rangle\langle0|\mathcal A_\kappa(\rho)]
   =\Tr[(1-\kappa)^H\rho]\le r/\kappa .
\]
Positivity of the truncated Kraus sum and the support definition yield
\(\Tr(Q\omega_\rho)\ge1-1/t_E-r/\kappa\).
Under reset, the vacuum term is orthogonal to \(Q\), and
\(\Tr[Q(\rho_R\otimes\tau_{\kappa,E,a})]
   \le a(K+1)r/(1+\kappa E)\).
The resulting branch decision error is bounded uniformly by
\[
 z_E=\lambda\left(\frac1{t_E}+\frac{\bar r_E}{\kappa_0}\right)
       +q\bar r_E\left(\frac{t_E}{\kappa_0}
                              +\frac2{1+\kappa_0 E}\right)\to0 .
                                                               \tag{38}
\]
For large \(E\), binary Fano bounds \(H(F|RB)\le h(z_E)\).
Writing \(p_*=1/(1+\kappa_0 E)\), (37)--(38) prove the all-input
capacity converse; thermal input gives the opposite bound:
\[
 f(a)-q h(p_*)-h(z_E)
 \le\lambda c_\kappa(E)-\CEA(\mathcal N_{\kappa,\tau_{\kappa,E,a}},H,E)
 \le f(a)+b(\lambda p_*,q).                                  \tag{39}
\]
It holds for all sufficiently large \(E\), uniformly in \(a,\kappa\).
The low-entropy inputs were treated before applying the finite-rank test.

\subsection{Uniform common codes and budget matching}
Use the input \(\sigma_E\) and window \(L\) of (5). Pure loss preserves
that finite window, so the fixed output decoding \(D_L\) acts exactly
on its transmitted branch. Positivity gives
\[
 \|\mathcal A_\kappa(\sigma_E)\|_\infty
 \le\frac1{[1+\kappa_0(E-1)](1-\epsilon)} .
\]
The three entropies in its mutual information differ from their thermal
values by a uniform amount tending to zero. Explicitly, the discarded
input tail has mass \(\epsilon\) and conditional mean \(E-1+L\);
source, receiver and complementary means are no larger than this.
Entropy of mixing bounds each entropy difference by
\(h(\epsilon)+\epsilon[g(E-1+L)+g(E-1)]\).
Changing mean \(E-1\) to \(E\) costs at most
\(2\log_2(1+1/(E-1))\) in \(c_\kappa\). Thus a valid uniform mutual
information error is
\(3\{h(\epsilon)+\epsilon[g(E-1+L)+g(E-1)]\}
     +2\log_2(1+1/(E-1))\to0\).
The sharper receiver branch bound (27) now applies with the displayed
transmitted operator norm. The rational/type construction of Section~3
works for the full finite channel index set \((\kappa,\tau)\), enforcing
the same physical input cost on every encoder output.

For Shannon budgets take \(m=\lfloor E/\ln(E+1)\rfloor\).
The matching witness weight is
\(a=\min\{1,\max\{0,(B-1)/g(\kappa E)\}\}\).
It is admissible by entropy of mixing, but need not lie below the
target \(t=\min\{1,B/\log_2(E+1)\}\).
With \(C_0=\max\{\log_2(1/\kappa_0),\log_2 e\}\), the absolute bound
\[
 |a-t|\le\min\left\{1,\frac1{g(\kappa_0 E)}
                              +\frac{C_0}{\log_2(E+1)}\right\}\to0
\]
is uniform in \(B,\kappa\). Use the modulus for \(f\) and the same
spectral tail as before. For fixed \(\alpha>1\), the exactly admissible
weight \(a=1-2^{-((\alpha-1)/\alpha)B}\) and \(m=\lfloor\sqrt E\rfloor\)
give the same matching, without a clipping error. Equation (39) and the
common-code lower bound prove (34).
Additional pure loss is receiver postprocessing, so \(c_\kappa\) increases
in \(\kappa\). Taking its infimum over the unknown parameter gives
\(c_{\kappa_0}\); a witness at \(\kappa_0\) matches that benchmark.
This proves (35).


\section{Limits, evidence and responsibility}
The zero-loss classification and the unassisted and two-way conclusions
concern the identity/replacement model. Section~\ref{sec:optics} extends
the assisted minimax curves to a quantum-limited attenuator/replacement
model. Its environment is vacuum and its transmissivity is bounded away
from zero. No thermal-loss or arbitrary memory channel is covered.
The vacuum--thermal resets are asymptotically least favorable for the
entanglement-assisted curves of Section~\ref{sec:curves}; no exact
finite-energy optimizer or uniqueness is proved. No nonzero-loss curve
is supplied for \(\alpha<1\). R\'enyi order and transmission
probability remain fixed. Efficient encoders, finite-block resource
estimates, universal two-way codes and peak-cost/local-energy constraints
are outside the conclusions.

The six supplementary programs, listed with their outputs in the
supplement's README, execute deterministic finite diagnostics in a writable
copy of their folder, using Python, NumPy, SciPy and mpmath.
They check interlacing,
noncommuting finite-code information bounds, spectral compression,
type-subspace information comparisons. The R\'enyi geometric/asymptotic
rows are evaluations of the displayed closed form, not independent
geometric-series equality assertions.
The additional minimax controls check noncommuting finite-support inputs
with exactly summed thermal tails; the scalar bound controls use 220-digit
arithmetic to resolve their largest energy scale. They do not optimize an
operational capacity or construct a universal code. The optical diagnostics
add finite-dimensional dilation, entropy-difference, branch-test and
thermal-truncation checks; their signal-array entries are bound slacks,
not capacity estimates.
They are corroboration, not proofs of an infinite-dimensional coding
limit. The analytical proof uses the explicitly cited coding and
distillation dependencies. Literature comparisons were directed and
selective, with no certification of worldwide priority.

\section*{Model assessments and review context}
\paragraph{AI and responsibility disclosure.}
The named author directed substantial OpenAI Codex assistance in derivation,
writing and numerical work. Requested GPT-6 Astra xhigh preparatory audits
examined the spectral/R\'enyi derivations and, in a separate subsequent
check, the compound-code extension, using shared files and reused context.
Those working-text audits were unscored and did not assess this combined
PDF. Inputs and public search results exposed historical evaluations of
other papers; no earlier score is adopted here.
No runtime isolation, independently attested model identity, independent
human review, Lean formalization, certified worldwide originality or
predicted assessment score is claimed. A full requested-model assessment
of the preceding frozen draft required a local correction to the
dephasing example. This revision makes that correction and clarifies
secrecy normalization, dependencies and the numerical evidence scope.
A further unscored working-argument audit examined the full minimax
loss derivation now included in Section~\ref{sec:curves}, with reused
context and shared files. It was not a full assessment of this PDF.
Another unscored working-argument audit examined the optical extension
of Section~\ref{sec:optics}, including the finite-rank branch test and
uniform transmissivity and budget remainders. It reused context and
shared files, and was not a full assessment of this combined PDF.
The preceding assessment and its objections are retained; its scores
are not transferred to this revised PDF. No full assessment of this exact
revised PDF has yet been obtained.

\begin{thebibliography}{99}
\small
\setlength{\itemsep}{1.5pt}
\bibitem{prior} L. Eriksson, \emph{Constant-gap entanglement distillation and
a spectral criterion for noise-universal communication},
AIRR.SCIENCE ARR-2026-58XZ9SCZA39DSRF0, v1 (2026), Working paper.
\url{https://airr.science/papers/ARR-2026-58XZ9SCZA39DSRF0/versions/v1/}.
\bibitem{minimax} L. Eriksson, \emph{Sharp minimax losses for energy-constrained
quantum reset channels}, AIRR.SCIENCE ARR-2026-2EXHFGWW0S8MG8QA,
v1 (2026), Working paper.
\url{https://airr.science/papers/ARR-2026-2EXHFGWW0S8MG8QA/versions/v1/}.
\bibitem{vc} F. Verstraete and J. I. Cirac, \emph{Matrix product states
represent ground states faithfully}, \href{https://arxiv.org/abs/cond-mat/0505140v6}
{arXiv:cond-mat/0505140v6} (2006).
\bibitem{schuch} N. Schuch, M. M. Wolf, F. Verstraete and J. I. Cirac,
\emph{Entropy scaling and simulability by Matrix Product States},
\href{https://arxiv.org/abs/0705.0292v2}{arXiv:0705.0292v2} (2008).
\bibitem{cq} I. Bjelakovi\'c and H. Boche,
\emph{Classical Capacities of Compound and Averaged Quantum Channels},
\href{https://arxiv.org/abs/0710.3027v2}{arXiv:0710.3027v2} (2009).
\bibitem{quantumcompound} I. Bjelakovi\'c, H. Boche and J. N\"otzel,
\emph{Entanglement transmission and generation under channel uncertainty:
Universal quantum channel coding},
\href{https://arxiv.org/abs/0811.4588v3}{arXiv:0811.4588v3} (2009).
\bibitem{eacompound} M. Berta, H. Gharibyan and M. Walter,
\emph{Entanglement-assisted capacities of compound quantum channels},
\href{https://arxiv.org/abs/1603.02282v2}{arXiv:1603.02282v2} (2017).
\bibitem{winter} A. Winter, \emph{Tight uniform continuity bounds for quantum
entropies: conditional entropy, relative entropy distance and energy constraints},
\href{https://arxiv.org/abs/1507.07775v6}{arXiv:1507.07775v6} (2016).
\bibitem{hs} A. S. Holevo and M. E. Shirokov,
\emph{On the entanglement-assisted classical capacity of infinite-dimensional
quantum channels}, \href{https://arxiv.org/abs/1210.6926}{arXiv:1210.6926} (2012).
\bibitem{dsw} N. Davis, M. E. Shirokov and M. M. Wilde,
\emph{Energy-constrained two-way assisted private and quantum capacities of
quantum channels}, \href{https://arxiv.org/abs/1801.08102v3}{arXiv:1801.08102v3}
(2018).
\bibitem{reliability} K. Watanabe, M. Hayashi, T. Matsuura and H.-C. Cheng,
\emph{Universality Sacrifices Reliability in Classical-Quantum Channel Coding},
\href{https://arxiv.org/abs/2610.01941v1}{arXiv:2610.01941v1} (2026).
\bibitem{glr} V. Giovannetti, S. Lloyd and M. B. Ruskai,
\emph{Conditions for the multiplicativity of maximal $\ell_p$-norms of
channels for fixed integer $p$},
\href{https://arxiv.org/abs/quant-ph/0408103v3}{arXiv:quant-ph/0408103v3} (2005).
\bibitem{cnr} A. Cacioppo, J. N\"otzel and M. Rosati,
\emph{Compound Channel Capacities under Energy Constraints and Application},
\href{https://arxiv.org/abs/2105.04274v1}{arXiv:2105.04274v1} (2021).
\bibitem{catalano} G. Catalano, M. Fanizza, F. A. Mele, G. De Palma and
V. Giovannetti, \emph{Convex combinations of bosonic pure-loss channels},
\href{https://arxiv.org/abs/2604.26874v2}{arXiv:2604.26874v2} (2026).
\bibitem{hw} A. S. Holevo and R. F. Werner,
\emph{Evaluating capacities of Bosonic Gaussian channels},
\href{https://arxiv.org/abs/quant-ph/9912067v1}
{arXiv:quant-ph/9912067v1} (1999), Sections IV.C and V.A.
\bibitem{gain} A. S. Holevo,
\emph{The entropy gain of infinite-dimensional quantum channels},
\href{https://arxiv.org/abs/1003.5765v1}{arXiv:1003.5765v1} (2010).
\end{thebibliography}
\end{document}
